\section{Recapitulation}

Saint Paul teaches in the epistle to the Colossians that all things will be summed up in Christ, who will then deliver the kingdom up to the Father, as He had it from Him from the beginning. Then, will come and be the ``end''. This teaching, along with verses such as ``heaven and earth will pass away, but the word of the Lord (My Word) shall endure forever'' give glimpses of another doctrine hiding behind the ``perspicacious'' meaning of \emph{sola Scriptura}. If ``heaven and earth'' will pass away, then there is an element of nihilism to Scripture, compatible with the ontological assertions of certain other ancient scriptures. Yet, the passing of heaven and earth is not ``the last word'' -- we can connect this with Tomberg's doctrine of the ``Self (God) beyond the ultimate Selves''. Not only can men like Enoch escape the horizontal circle of the serpent, so can \emph{buddhas} from the East, and yet their escape remains an inexplicable fact. Avatars are introduced to communicate in the opposite direction, but this, too, is a mystery. Men, said Paul on Mars Hill\footnote{\url{https://en.wikipedia.org/wiki/Areopagus}}, quoting Aratus and other Greek poets, seek for God, if ``happily they might find him''. The possibility is real, and not limited to the Hebraic tradition. However, the esoteric Christian tradition is not so much that others (such as the Greek) are ``incomplete'', but that it is capable of being added to. \emph{These are not quite the same things, to be incomplete and to be capable of addition.}

But this is not all. From another Pauline epistle comes this\footnote{\url{http://biblebrowser.com/1\_corinthians/3-12.htm}}: ``man's works shall be burned up in fire, but he shall be saved''.

The \emph{Apokastasis}\footnote{\url{https://en.wikipedia.org/wiki/Apocatastasis}} used to be the doctrine of the early church, until the needs of holding together an Empire and a marching order of transcendent temple-worship lead to Justinian's measures\footnote{\url{http://en.wikipedia.org/wiki/Justinian\_I}}. The West was in such dire shape following the barbarian tribal invasions and the imperial collapse that Church fasting requirements were relaxed so that the sick and the poor could eat enough meat to survive -- this was common in the West, at least. This was the kind of context into which Christianity stepped to carry the banner of transcendence and order. The original vision of a recapitulation (or re-breathing of God's breath) gave place to more concrete, useful, and also (in their own way) more accurate reflections of what man needed to know during those dire times. What good was a doctrine of universal salvation if it would be simply used to justify laxity and apathy?

Scripture teaches that Christ loved his sheep (His ``body'') as a bridegroom loves his bride-to-be. If Christ intended to not enter heaven without cleansing his bride (and even ``harrowing hell''\footnote{\url{https://en.wikipedia.org/wiki/Harrowing\_of\_Hell}}), this leads us to believe that the ``path'' which he opened as an avatar was intended to embrace the entirety of the human race (including other avatars with a human face), because ``if I shall be lifted up, I shall draw all men to me''.

The foregoing overview should make clear the complexity and difficulty of the Christian tradition. It is not so simple as to say that the Church did merely or only this, or that, or that Christ taught only this, or just that. There were multiple facets or aspects of the primordial Tradition being preserved, but being preserved while undergoing a metamorphosis. If we object that much was lost, indeed this is true, yet how can it be ``lost'' if we know it was? Esoteric traditions, such as the Fedeli Amore\footnote{\url{http://web.eecs.utk.edu/~mclennan/Classes/US310/Dante-Fedeli-d-Amore.html}}, have persisted within Christendom throughout history, down to this day, and one can trace its influence (as above) in the writings of Saint Paul, who had ``many things hard to understand'' and many other things he ``wished to say to you'', but ``you were not able''.

Thus, the nihilism of the ``heaven and earth being rolled up as a scroll'' is consistent with an esoteric doctrine of Christ as final avatar, not to abolish other avatars, but to entirely make sense of them and complete them by adding the cornerstone, which is as great a ``mystery beyond the mystery'' as esoteric doctrine is a ``mystery'' to something purely external.

In fact, the doctrine of Recapitulation can be deduced from that of karma and sexual union. If man and woman ``become one flesh'' (as Saint Paul warns those thinking of undertaking temple prostitution), then the subtle and vital bodies (as well as the physical) are united. Christianity, in fact, does not deny the ancient yoga techniques of retention of semen (the physical body), holding of the breath (the subtle body), and the stopping of thought (the vital body). In fact, it incorporated it into its priesthood. But it argued that there was ``another way''. This is true Christian form. As long as the reason is understood all is well. And here it is. If a man loves his woman as Christ loved the Church (with the result that the ``one flesh'' is totally sanctified, a double burden, or perhaps a double help, in the end), then the karmic chain of carnal relations to the woman will not hold the male out of paradise. Just as Christ ``weds'' to the Church, so can Christians make of monogamy a chivalric chastity and ``second coming of Love'', which can inspire something akin to that which fired the ``Dark Ages'' of Dante.

So which path is easier? Total abstinence, or devotion to ``the one woman'' above all else? Either way, it is time for Christians to repent.


\flrightit{Posted on 2012-03-03 by Logres}

\begin{center}* * *\end{center}

\begin{footnotesize}
\begin{sffamily}
\texttt{Logres on 2012-03-04 at 00:49 said:}

Much more really needs to be said and meditated upon, and this is (in a sense) woefully incomplete. For one thing, Mouravieff has interesting passages in Gnosis on the fact that ``a man is not without the woman, nor the woman without the man, in the Lord''. HOO might add something here about Percival, perhaps. Von Hildebrand (a convert from Protestantism to Catholicism) has interesting meditations on the man-woman dyad, and Steiner himself remarked that the male spiritual body was feminine, while the woman's was male. What needs more meditation is the fact that Christ did not come as other avatars did, withdrawing into himself upon a mountaintop or waging war with magic upon the kingdoms of the world (although he implicitly endorses such, in its place), but as an avatar of avatar who incorporated a multitude of previous faces/masks and made a particular Tao to unify what had before been at an impasse. In particular, the bride-bridegroom motif (which is visible in historical form in the history ofAnglicanism in Britain, most embodied, perhaps in the writings of Jane Austen) seems to be a singularity which ought only to be attempted by God.

\hfill

\end{sffamily}
\end{footnotesize}

